\begin{figure}[htbp]
\centering
\begin{tikzpicture}[
  x=35cm,y=.86cm,font=\small,
  slot/.style={black,line width=.87pt},
  guide/.style={black!48,densely dashed,line width=.4pt},
  axis/.style={black!55,line width=.45pt},
  bracket/.style={line width=.8pt}]
 % Show the first .35 of the unit mass interval at its actual scale.
 \fill[black!13] (.1875,0) rectangle (.20,3.38);
 \draw[guide] (0,1.85) -- (.35,1.85) node[right,text=black] {\(t'\)};
 \draw[guide] (0,1.12) -- (.35,1.12) node[right,text=black] {\(t\)};
 \node[anchor=east] at (0,3.05) {\(d_s(y)\)};
 \foreach \a/\b/\h in {
   0/.04/3.05,.04/.095/2.75,.095/.15/2.42,
   .15/.19375/1.85,.19375/.25/1.85,.25/.285/1.46,
   .285/.325/.78,.325/.35/.50}
   \draw[slot] (\a,\h) -- (\b,\h);
 % A vertical tick separates the two slots with the same surplus.
 \draw[slot] (.19375,1.74) -- (.19375,1.96);
 \fill (.1895,1.85) circle[radius=1.5pt];

 \draw[axis] (0,0) -- (.35,0);
 \foreach \x/\lab in {0/0,.10/0.10,.1875/\gamma,.20/\delta,.30/0.30,.35/0.35}{
   \draw[axis] (\x,0) -- (\x,-.1);
   \node[below=2pt,inner sep=1pt] at (\x,0) {\(\lab\)};
 }
 \node[anchor=north,font=\footnotesize] at (.296,-.44) {Cumulative slot mass};
 \draw[bracket,lagoon] (0,3.52) -- (0,3.68) -- (.25,3.68) -- (.25,3.52);
 \node[above=1pt,text=lagoon] at (.125,3.68) {\(v_s=1\) if the child succeeds};
 \draw[bracket,ember] (0,-.52) -- (0,-.68) -- (.20,-.68) -- (.20,-.52);
 \node[below=1pt,text=ember] at (.10,-.68) {Multiply this mass by four};
\end{tikzpicture}
\caption{A heavy probe, with the first \(0.35\) of the slot mass shown.
Each horizontal segment represents a slot, ordered by surplus. The dot selects
a pivot in the shaded band \([\gamma,\delta)=[3/16,1/5)\).
Two adjacent slots have the pivot surplus \(t'\). Success at \(t'\) guarantees
correct feedback on mass \(0.25\), including both tied slots.
Reweighting multiplies the first \(0.20\) of the mass by four before normalization.}
\label{fig:band}
\end{figure}
